\chapter{G\'en\'eration augment\'ee par r\'ecup\'eration (RAG)}
\label{ch:rag}

\begin{quote}
\emph{<< Le RAG permet \`a un LLM de r\'epondre depuis vos documents au lieu d'inventer. C'est le pattern d'IA g\'en\'erative le plus pratique en production aujourd'hui. >>}
\end{quote}

% --- hook ---
Un LLM ne connaît rien après sa date d'entraînement. Posez-lui une question sur vos données internes, et il \emph{inventera} une réponse plausible. Le RAG résout les deux problèmes en une étape : on récupère d'abord les documents pertinents, puis on génère la réponse en s'appuyant dessus. C'est aujourd'hui le pattern d'IA générative le plus déployé en production --- bien plus que le fine-tuning, parce qu'il est plus simple, plus traçable, et plus facile à mettre à jour.

% ============================================================
\section{Pourquoi le RAG ?}

Les LLM ont deux limites fondamentales :
\begin{enumerate}
  \item \textbf{Date de coupure de connaissance.} Ils ne savent rien apr\`es leur date d'entra\^inement.
  \item \textbf{Hallucination.} Ils g\'en\`erent avec assurance des informations plausibles mais fausses.
\end{enumerate}

Le RAG r\'esout les deux : r\'ecup\'erer des documents pertinents, puis g\'en\'erer une r\'eponse ancr\'ee dans ces documents.

\begin{center}
\begin{tikzpicture}[
  block/.style={rectangle, draw, fill=aiblue!15, minimum width=2.8cm, minimum height=1cm, rounded corners, align=center},
  arrow/.style={-{Stealth}, thick}
]
  \node[block] (query) at (0,0) {Requ\^ete\\utilisateur};
  \node[block] (embed) at (3.5,0) {Plonger\\la requ\^ete};
  \node[block] (search) at (7,0) {Recherche\\vectorielle};
  \node[block] (docs) at (7,-2) {Fragments\\de documents};
  \node[block] (prompt) at (10.5,0) {Construire\\le prompt};
  \node[block] (llm) at (10.5,-2) {LLM\\g\'en\`ere};

  \draw[arrow] (query) -- (embed);
  \draw[arrow] (embed) -- (search);
  \draw[arrow] (search) -- (prompt);
  \draw[arrow] (docs) -- (search);
  \draw[arrow] (prompt) -- (llm);
\end{tikzpicture}
\end{center}

% ============================================================
\section{Plongements pour la r\'ecup\'eration}

On utilise \textbf{sentence-transformers} pour convertir le texte en vecteurs denses. Des textes similaires ont des vecteurs similaires (haute similarit\'e cosinus).

\begin{minted}{python}
from sentence_transformers import SentenceTransformer
import numpy as np

model = SentenceTransformer("all-MiniLM-L6-v2")

sentences = [
    "The patient has a high fever and cough.",
    "Machine learning is a subset of artificial intelligence.",
    "The patient presents with elevated temperature and respiratory symptoms.",
]

embeddings = model.encode(sentences)
print(f"Embedding shape: {embeddings.shape}")  # (3, 384)

# Similarit\'e cosinus
from sklearn.metrics.pairwise import cosine_similarity
sim_matrix = cosine_similarity(embeddings)
print("Similarity matrix:")
print(np.round(sim_matrix, 3))
# les phrases 0 et 2 doivent avoir une similarit\'e \'elev\'ee
\end{minted}

\begin{aitip}
\texttt{all-MiniLM-L6-v2} est un mod\`ele de plongement l\'eger (80 Mo) qui tourne vite sur CPU. Pour plus de qualit\'e, consid\'erez \texttt{BAAI/bge-small-en-v1.5} ou \texttt{nomic-ai/nomic-embed-text-v1.5}.
\end{aitip}

% ============================================================
\section{Bases de donn\'ees vectorielles : ChromaDB}

ChromaDB est une base vectorielle open-source qui tourne localement sans configuration :

\begin{minted}{python}
import chromadb
from chromadb.utils import embedding_functions

# Cr\'eer un client persistant
client = chromadb.PersistentClient(path="./chroma_db")

# Utiliser sentence-transformers pour les plongements
ef = embedding_functions.SentenceTransformerEmbeddingFunction(
    model_name="all-MiniLM-L6-v2"
)

# Cr\'eer une collection
collection = client.get_or_create_collection(
    name="course_notes",
    embedding_function=ef,
)

# Ajouter des documents
documents = [
    "The Transformer architecture uses self-attention mechanisms.",
    "LoRA reduces fine-tuning costs by training low-rank adapters.",
    "RAG retrieves relevant documents before generating answers.",
    "Diffusion models generate images by reversing a noise process.",
    "RLHF aligns language models with human preferences.",
]

collection.add(
    documents=documents,
    ids=[f"doc_{i}" for i in range(len(documents))],
    metadatas=[{"chapter": i+1} for i in range(len(documents))],
)

# Requ\^ete
results = collection.query(query_texts=["How does fine-tuning work?"], n_results=2)
print("Top results:")
for doc, dist in zip(results["documents"][0], results["distances"][0]):
    print(f"  [{dist:.4f}] {doc}")
\end{minted}

% ============================================================
\section{Bases de donn\'ees vectorielles : FAISS}

FAISS (Facebook AI Similarity Search) est optimis\'e pour la recherche de similarit\'e \`a grande \'echelle :

\begin{minted}{python}
import faiss
import numpy as np
from sentence_transformers import SentenceTransformer

model = SentenceTransformer("all-MiniLM-L6-v2")
documents = [
    "Python is a high-level programming language.",
    "PyTorch is a deep learning framework.",
    "LangChain helps build LLM applications.",
    "ChromaDB stores vector embeddings.",
    "Transformers use attention mechanisms.",
]

# Encoder et construire l'index
embeddings = model.encode(documents).astype("float32")
dimension = embeddings.shape[1]

index = faiss.IndexFlatL2(dimension)  # distance L2 (euclidienne)
index.add(embeddings)
print(f"Index contains {index.ntotal} vectors")

# Recherche
query = model.encode(["How do I build an LLM app?"]).astype("float32")
distances, indices = index.search(query, k=2)
for idx, dist in zip(indices[0], distances[0]):
    print(f"  [{dist:.4f}] {documents[idx]}")
\end{minted}

% ============================================================
\section{D\'ecoupage de documents}

Les vrais documents sont trop longs pour un seul plongement. On les d\'ecoupe en fragments qui se chevauchent :

\begin{minted}{python}
from langchain_text_splitters import RecursiveCharacterTextSplitter

text = open("my_document.txt").read()  # ou n'importe quel long texte

splitter = RecursiveCharacterTextSplitter(
    chunk_size=500,       # caract\`eres par fragment
    chunk_overlap=50,     # chevauchement entre fragments
    separators=["\n\n", "\n", ". ", " ", ""],
)

chunks = splitter.split_text(text)
print(f"Number of chunks: {len(chunks)}")
print(f"First chunk ({len(chunks[0])} chars):\n{chunks[0][:200]}...")
\end{minted}

\begin{warningbox}
La taille des fragments est critique. Trop petite : contexte perdu. Trop grande : pertinence dilu\'ee. Commencez \`a 500--1000 caract\`eres et ajustez selon votre qualit\'e de r\'ecup\'eration.
\end{warningbox}

% ============================================================
\section{Pipeline RAG complet avec LangChain}

\begin{minted}{python}
# !pip install langchain langchain-community langchain-google-genai chromadb

from langchain_community.document_loaders import PyPDFLoader
from langchain_text_splitters import RecursiveCharacterTextSplitter
from langchain_community.vectorstores import Chroma
from langchain_community.embeddings import HuggingFaceEmbeddings
from langchain_google_genai import ChatGoogleGenerativeAI
from langchain_core.prompts import ChatPromptTemplate
from langchain_core.runnables import RunnablePassthrough
from langchain_core.output_parsers import StrOutputParser

# 1. Charger le PDF
loader = PyPDFLoader("research_paper.pdf")
pages = loader.load()
print(f"Loaded {len(pages)} pages")

# 2. D\'ecouper en fragments
splitter = RecursiveCharacterTextSplitter(chunk_size=800, chunk_overlap=100)
chunks = splitter.split_documents(pages)
print(f"Created {len(chunks)} chunks")

# 3. Cr\'eer le vector store
embeddings = HuggingFaceEmbeddings(model_name="all-MiniLM-L6-v2")
vectorstore = Chroma.from_documents(chunks, embeddings, persist_directory="./rag_db")
retriever = vectorstore.as_retriever(search_kwargs={"k": 4})

# 4. Cr\'eer la cha\^ine RAG
llm = ChatGoogleGenerativeAI(model="gemini-2.0-flash", temperature=0.3)

template = ChatPromptTemplate.from_template("""Answer the question based only on the following context. If you cannot answer from the context, say "I don't have enough information."

Context:
{context}

Question: {question}

Answer:""")

def format_docs(docs):
    return "\n\n".join(doc.page_content for doc in docs)

rag_chain = (
    {"context": retriever | format_docs, "question": RunnablePassthrough()}
    | template
    | llm
    | StrOutputParser()
)

# 5. Interroger
answer = rag_chain.invoke("What is the main contribution of the paper?")
print(answer)
\end{minted}

% ============================================================
\section{\'Evaluer la qualit\'e d'un RAG}

\begin{minted}{python}
# \'Evaluation simple de la r\'ecup\'eration : v\'erifier que le bon document est trouv\'e
test_questions = [
    {"question": "What is LoRA?", "expected_doc_id": 1},
    {"question": "How does RAG work?", "expected_doc_id": 2},
]

correct = 0
for test in test_questions:
    results = collection.query(query_texts=[test["question"]], n_results=1)
    retrieved_id = int(results["ids"][0][0].split("_")[1])
    if retrieved_id == test["expected_doc_id"]:
        correct += 1

print(f"Retrieval accuracy: {correct}/{len(test_questions)}")
\end{minted}

\begin{exercise}
\begin{enumerate}
  \item Construisez une application RAG sur 3 fichiers PDF de votre choix. Utilisez ChromaDB pour le stockage et Gemini pour la g\'en\'eration. Testez avec 5 questions.
  \item Comparez la qualit\'e de r\'ecup\'eration entre \texttt{all-MiniLM-L6-v2} et \texttt{BAAI/bge-small-en-v1.5} sur le m\^eme corpus.
  \item Exp\'erimentez avec des tailles de fragments (200, 500, 1000, 2000 caract\`eres). Mesurez la pr\'ecision de r\'ecup\'eration pour 10 questions de test \`a chaque taille.
  \item Ajoutez un filtrage par m\'etadonn\'ees \`a votre cha\^ine RAG : r\'ecup\'erer seulement depuis des chapitres ou types de documents sp\'ecifiques.
  \item Impl\'ementez une fonctionnalit\'e << sources >> : apr\`es g\'en\'eration, montrer \`a l'utilisateur quels fragments ont \'et\'e utilis\'es.
\end{enumerate}
\end{exercise}

\begin{ethicsbox}
Le RAG r\'eduit l'hallucination mais ne l'\'elimine pas. Le LLM peut toujours mal interpr\'eter le contexte r\'ecup\'er\'e, combiner l'information incorrectement, ou ajouter de l'information absente des documents. Fournissez toujours des citations de sources pour que les utilisateurs puissent v\'erifier. Dans les domaines \`a fort enjeu (m\'edical, juridique, financier), la sortie RAG doit \^etre revue par un humain qualifi\'e.
\end{ethicsbox}

% ============================================================
\section{R\'esum\'e du chapitre}

\begin{itemize}
  \item Le RAG combine r\'ecup\'eration (trouver les documents pertinents) et g\'en\'eration (r\'epondre sur la base de ces documents).
  \item Sentence-transformers convertit le texte en vecteurs denses pour la recherche s\'emantique.
  \item ChromaDB et FAISS sont des bases vectorielles pour stocker et chercher des plongements.
  \item Le d\'ecoupage de documents segmente les longs textes en fragments qui se chevauchent.
  \item LangChain fournit un pipeline RAG complet : charger, d\'ecouper, plonger, stocker, r\'ecup\'erer, g\'en\'erer.
  \item Taille de fragment, qualit\'e du mod\`ele de plongement et nombre de r\'esultats r\'ecup\'er\'es ($k$) sont les param\`etres cl\'es \`a r\'egler.
  \item Le RAG r\'eduit l'hallucination mais demande des citations de sources pour la fiabilit\'e.
\end{itemize}
